\documentclass[10.5pt,letterpaper]{article}

\usepackage[margin=0.72in]{geometry}
\usepackage[T1]{fontenc}
\usepackage[utf8]{inputenc}
\usepackage{lmodern}
\usepackage[hidelinks]{hyperref}
\usepackage{enumitem}
\usepackage{titlesec}
\usepackage{xcolor}

\definecolor{sectiongray}{gray}{0.20}
\setlength{\parindent}{0pt}
\setlength{\parskip}{4pt}
\setlist[itemize]{leftmargin=1.20em, topsep=2pt, itemsep=1pt, parsep=0pt}
\titleformat{\section}{\large\bfseries\color{sectiongray}}{}{0em}{}[\titlerule]
\titleformat{\subsection}{\normalsize\bfseries}{}{0em}{}
\titlespacing*{\section}{0pt}{6pt}{3pt}
\titlespacing*{\subsection}{0pt}{5pt}{1pt}

\newcommand{\project}[4]{%
\subsection*{#1 \hfill {\small #2}}
\textbf{Current output:} #3\\
\textbf{Near-term needs:} #4
}

\begin{document}

{\LARGE \textbf{Paul A. Skeffington --- Research Status}}\\[-1pt]
{\small Public-facing research storyboard, current projects, publication status, and near-term needs register.}\\[6pt]

\section*{Research Storyboard}
This portfolio is organized around careful, reproducible use of public-health, infrastructure, environmental, administrative, longitudinal, biomedical, security-evaluation, and public-history data. The current emphasis is on building transparent workflows, documenting evidence boundaries, and moving selected projects from exploratory scaffolds toward validated manuscripts, methods notes, report artifacts, or decision-support tools.

\section*{Current Project Board}
\project{Life-Course Data, Family Economics, and Fertility}{Manuscript cleanup}{NLSY79 manuscript workspace with fixed-effects, count-model, survival-analysis, LaTeX, and table workflow.}{Citation verification; table-to-model audit; journal-format polish; reviewer-response planning.}

\project{Rural Water, Wastewater, and Infrastructure Health Risk}{Active / pending validation}{McDowell County GIS and decision-support workflow for terrain/access scoring, stream-aware wastewater indices, stakeholder planning, and source-validation tracking.}{Official water-quality confirmation; utility/source-holder records; final variable-confidence ledger.}

\project{Humanitarian WASH and Health-System Disruption}{Presentation complete / methods-note candidate}{Gaza WASH rapid assessment framework using public humanitarian, WASH, displacement, environmental-health, and geospatial sources.}{Dated source manifest; limitations section; separation of indexed inputs from interpretive claims.}

\project{Haiti Nippes Public-Health and Infrastructure Assessment}{New public repository / intake needed}{Intake-stage public repository for regional public-health, infrastructure, and humanitarian assessment work.}{Complete source manifest, project boundary, claim controls, and intended public output.}

\project{Computational Diagnostics and Cipher Topology}{Active scaffold}{Cipher-topology lab using persistent-homology features as auxiliary diagnostics for structured or biased cipher-output controls.}{Freeze corpus; document null/control design; add notebooks; define contribution boundary.}

\project{Identity-Abuse and Misuse-Risk Evaluation}{New public repository / intake needed}{Intake-stage public repository for identity-abuse, misuse-risk, or safety-evaluation methods.}{Define public-safe threat model, evidence boundary, and executable validation target.}

\project{Medical-Signal Noise Reduction}{Early-stage topic selection}{ECG denoising / signal-quality direction aligned with open medical data and reproducible benchmark design.}{Select dataset; define noise classes; benchmark methods; create results table.}

\project{PET Noise and Radiomics Robustness}{Early-stage scaffold}{Biomedical imaging scaffold focused on simulated PET noise, lesion segmentation, and radiomics feature stability.}{Select public PET/CT dataset; define perturbation protocol; separate exploratory claims from validated results.}

\project{Cancer End-of-Life Typologies}{Early-stage topic selection}{Open-data scaffold aligned to place-of-death, cancer care, and end-of-life quality variation.}{Identify dataset; define typology outcome; map variables; avoid unsupported causal framing.}

\project{Public-Health Emergency Preparedness and Practicum Reporting}{Public report artifacts / intake needed}{Public-facing report and preparedness repositories for emergency-response, practicum, and public-health evidence packaging.}{Add source manifests, project boundaries, and public-safe deliverable targets.}

\project{Catholic Archive and Public-History Indexing}{Active archive build}{St. Bonaventure reliquary section mapping, visible-label transcription, entity modeling, crop tooling, and provenance controls.}{Complete section index; distinguish devotional tradition from verified history; add uncertainty flags.}

\section*{Selected Research Outputs}
\cvrole{CART-TRACE: longitudinal CAR T-cell care trajectories}{Final public synthetic scholarly freeze}
\cvplace{Reproducible translational health data-science capstone}

\cvrole{Family economics, financial literacy, and fertility outcomes using NLSY79 data}{Pre-submission review}
\cvplace{Longitudinal manuscript in blinded-package and citation review}

\cvrole{Humanitarian WASH, displacement, and health-system disruption}{Methods note in development}
\cvplace{Public-source humanitarian and environmental-health methods}

\cvrole{Rural water, wastewater, and infrastructure health risk}{Completed applied public-health practicum}
\cvplace{Public-safe project claim; scholarly extensions handled separately}

\cvrole{Biomedical data-science methods}{Developing under explicit evidence gates}
\cvplace{ECG benchmark pathway, PET imaging robustness, and cancer end-of-life methods}


\section*{Completed / Print-Ready Objects}
\begin{itemize}
    \item Public CV package with standard, one-page, index-safe, and research-status outputs.
    \item Recursive repository intake layer in the private CV workspace, with newly discovered public repositories carried forward conservatively as intake-stage objects.
    \item Gaza WASH presentation object suitable for conversion into a reproducible methods note.
    \item McDowell GIS and stakeholder scaffolds, pending official-source validation.
\end{itemize}

\section*{Pending Validation}
\begin{itemize}
    \item Water-quality and wastewater indicators remain secondary-extracted until confirmed through official or certified sources.
    \item NLSY79 manuscript claims require final citation, model-table, and sample-audit parity.
    \item Newly discovered public repositories remain intake/status objects until datasets, source manifests, methods, and claim boundaries are frozen.
\end{itemize}

\section*{Immediate Major Needs}
\begin{itemize}
    \item Wire recursive inventory output into the public export path so new public repositories are picked up before PDF builds.
    \item Complete project-intake records for newly discovered public repositories before adding strong CV claims.
    \item For McDowell, obtain official water-system and wastewater confirmation before using measured-risk language.
    \item For NLSY79, lock manuscript tables and preserve a direct audit trail from model output to LaTeX tables.
    \item For medical/noise projects, select public datasets before expanding claims.
\end{itemize}

\end{document}
